\begin{pastexam}[label={px:ntut-113-ee-ds-07}]{北科大電機所（戊組）113 學年度 第 II.3 題}
（簡答；4 分）

What is the time complexity of the dynamic programming implementation of the longest common
subsequence problem where the length of one string is ``m'' and the length of the other string is
``n'' ? (4\%)
\end{pastexam}

\begin{pastsoln}
答案是 $\Theta(mn)$。

設兩個字串是 $X = x_1 x_2 \cdots x_m$、$Y = y_1 y_2 \cdots y_n$，$c[i, j]$ 是 $X$ 的前 $i$ 個字元和
$Y$ 的前 $j$ 個字元的 LCS 長度（資結書第~18.3~節）：
\[
c[i, j] = \begin{cases}
0, & i = 0 \text{ 或 } j = 0, \\
c[i-1, j-1] + 1, & i, j > 0 \text{ 且 } x_i = y_j, \\
\max(c[i-1, j],\ c[i, j-1]), & i, j > 0 \text{ 且 } x_i \ne y_j.
\end{cases}
\]
表格有 $(m + 1)(n + 1)$ 格，每一格只看左、上、左上三格，花 $O(1)$。 所以時間是
$\Theta((m+1)(n+1)) = \Theta(mn)$；表格空間也是 $\Theta(mn)$。 如果只要長度、不必印出 LCS，可以只留
兩列，空間降到 $\Theta(\min(m, n))$，時間不變。 從 $c[m, n]$ 往回追出一個 LCS 另外花 $O(m + n)$。
\end{pastsoln}
